% ============================================================================
%  Anexos
% ============================================================================
% Cambiar aquí la dirección del repositorio
\newcommand{\repourl}{https://github.com/EstebanUnal-Hub/ADRC_WWTP}
\newcommand{\repositorio}{\url{\repourl}}
% enlace a una sección del README: \readme{ancla}{texto}
\newcommand{\readme}[2]{\href{\repourl\##1}{#2}}

\chapter{Repositorio de código y datos}
\label{anx:st}
\label{anx:archivos}

El código, los cuadernos, los datos procesados y los bloques del PLC de este trabajo están en el
repositorio \repositorio, cuyo \href{\repourl}{README} los explica en detalle. La
Tabla~\ref{tab:anx_repo} resume los archivos principales; la última columna enlaza a la sección del
README que explica cada uno.

{\renewcommand{\arraystretch}{1.15}
\begin{longtable}{|>{\footnotesize\raggedright\arraybackslash}p{5.6cm}|>{\small\raggedright\arraybackslash}p{6.9cm}|>{\small\centering\arraybackslash}p{1.9cm}|}
\caption{Archivos principales del repositorio.}
\label{tab:anx_repo} \\
\hline
\rowcolor[gray]{0.9}
\small\textbf{Archivo} & \textbf{Contenido} & \footnotesize\textbf{README} \\ \hline
\endfirsthead
\hline
\rowcolor[gray]{0.9}
\small\textbf{Archivo} & \textbf{Contenido} & \footnotesize\textbf{README} \\ \hline
\endhead
\texttt{README.md} & Descripción general del repositorio y de cada carpeta. & \href{\repourl}{Inicio} \\ \hline
\texttt{Datos/analisis/}\newline\texttt{analisis\_dinamica\_limpio.ipynb} & Identificación de la planta
por ventanas ARX y caracterización con el histórico del SCADA (Capítulo~\ref{cap:caracterizacion}). &
\readme{3-plant-model-and-identification}{Sec. 3} \\ \hline
\texttt{Simulaciones/PTAR/02 Diseño/Python/}\newline\texttt{ADRC\_PTAR\_continuo.ipynb} & ADRC con ESO
y GPI en tiempo continuo: diseño, equivalencia, \emph{anti-windup}, acondicionamiento de la referencia
y comparación con el PID (Capítulos~\ref{cap:ADRC-ESO} a \ref{cap:comparacion_pid}). &
\readme{4-error-based-adrc-with-extended-state-observer}{Sec. 4 a 6} \\ \hline
\texttt{Simulaciones/PTAR/02 Diseño/Python/}\newline\texttt{ADRC\_PTAR\_discreto.ipynb} & ADRC con ESO
y GPI en tiempo discreto: discretización y simulación con retardo
(Sección~\ref{sec:impl_discretizacion}). & \readme{64-discretization}{Sec. 6.4} \\ \hline
\texttt{P\&ID/P\&ID.vsdx} y \texttt{P\&ID.pdf} & Plano de la etapa de aireación, con el lazo 100 de
oxígeno disuelto (Figura~\ref{fig:pid_tikz}). & \readme{2-process-and-control-loop}{Sec. 2} \\ \hline
\texttt{PLC\_ST/Codigos/*.st}\newline\texttt{PLC\_ST/Libreria\_ADRC/} & Los siete bloques de función en
texto estructurado (IEC~61131-3), también documentados como librería de ctrlX
(Sección~\ref{sec:impl_plc}). & \readme{7-plc-implementation}{Sec. 7} \\ \hline
\end{longtable}
}
